\documentclass[10pt,a4paper]{amsart}
\usepackage[margin=19mm]{geometry}
\usepackage{amsmath,amssymb,mathtools}
\usepackage[scr=rsfs,cal=euler]{mathalfa}
\usepackage{xfrac}
\usepackage[unicode,hidelinks,bookmarks=false]{hyperref}
\newcommand{\ie}{i.e.\ }
\newcommand{\cf}{cf.\ }
\newcommand{\resp}{respectively\ }
\newcommand{\Subsection}{\subsection}
\providecommand{\nobreakdash}{\nobreak\hbox{-}}
\setlength{\emergencystretch}{2em}
\multlinegap0pt
\usepackage{amssymb}
\usepackage{bm}
\usepackage{xcolor}
\usepackage[shortlabels]{enumitem}
\usepackage{tikz}
\newtheorem{lemma}{Lemma}
\newtheorem{claim}{Claim}
\newtheorem{midproof}{Midproof}
\renewcommand{\c}[1]{\langle #1 \rangle}
\newcommand{\func}{\operatorname}
\newcommand{\supp}{\func{supp}}
\title{Full mad families of vector spaces and two local Ramsey theories}
\author{Clement Yung}
\date{Source: arXiv:2503.24207v5 (CC BY 4.0)}
\begin{document}
\maketitle

\noindent\emph{Source and licence.} Full mad families of vector spaces and two local Ramsey theories. Version arXiv:2503.24207v5 (CC BY 4.0), distributed by arXiv under the Creative Commons Attribution 4.0 International licence, http://creativecommons.org/licenses/by/4.0/, as recorded in the arXiv licence metadata for this exact version and shown on the abstract page https://arxiv.org/abs/2503.24207v5. Full licence notice: \texttt{licenses/arxiv-2503.24207-CC-BY-4.0.txt}.\par\smallskip

\noindent\textit{Editorial scope.} Counted unit: the article's lemma lem:block.subspaces.intersection.lemma (source0.tex lines 410--413). The article's complete original proof of it is delivered with it (source0.tex lines 415--549, one \texttt{proof} environment of 135 lines). No mathematics is written, repaired or re-derived: the statement and the proof are the article's own lines, in the article's own notation, and no retained line is substituted or re-flowed.  No further source-local context is needed: every cross-reference of every retained line resolves inside the two delivered spans.  Nothing was omitted from any retained span, so the package carries no omission record.  No retained line contains a citation, so the wrapper carries no bibliography and no bibliography entry.  No retained line contains a figure environment, an image-inclusion command or drawing code, so no image is delivered and the package declares no asset dependency.  Editorial formatting only, in addition: an AMS \texttt{amsart} preamble that restates the article's own declarations and macros used by the retained lines, namely the environments \texttt{lemma}, \texttt{claim}, \texttt{midproof} on separate counters, together with the standard packages those lines need.  The retained environments print in this document as Lemma 1, Claim 1, Midproof 1 in order of appearance, whereas the article prints the same items with the numbering of its own full document.  The complete native source of the article as distributed in the arXiv e-print for this exact version is bundled unchanged as \texttt{sources/175701/source0.tex}.
\begin{lemma}
\label{lem:block.subspaces.intersection.lemma}
    Let $A,B \in E^{[\infty]}$ be such that $\c{A} \cap \c{B}$ is infinite-dimensional. Then there exists some $C \in E^{[\infty]}$ such that $\c{C} = \c{A} \cap \c{B}$.
\end{lemma}

\begin{proof}
    We write $A = (x_n)_{n<\omega}$ and $B = (y_n)_{n<\omega}$. For each $v \in \c{A} \cap \c{B}$, we let:
    \begin{align*}
        \supp_A(v) &:= \{n < \omega : \supp(x_n) \subseteq \supp(v)\}, \\
        \supp_B(v) &:= \{n < \omega : \supp(y_n) \subseteq \supp(v)\}.
    \end{align*}
    We remark that for all $v,w \in \c{A} \cap \c{B}$:
    \begin{align*}
        \supp(v) \subseteq \supp(w) &\iff \supp_A(v) \subseteq \supp_A(w) \\
        &\iff \supp_B(v) \subseteq \supp_B(w).
    \end{align*}
    For an arbitrary $v \in E$ and $k \in \supp(v)$, recall that $\lambda_k(v) \in \mathbb{F}$ denotes the unique coefficient such that $v = u + \lambda_k(v)e_k + w$ for some vectors $u < e_k < w$ (or $u = 0$ or $w = 0$).

    \begin{claim}
    \label{claim:supp(v).intersection}
        Let $v_0,v_1 \in \c{A} \cap \c{B}$, and suppose that $\supp_A(v_0) \cap \supp_A(v_1) \neq \emptyset$. Then there exists some $w \in \c{A} \cap \c{B}$ such that $\supp_A(w) = \supp_A(v_0) \cap \supp_A(v_1)$.
    \end{claim}

    \begin{midproof}
        Let $v_0,v_1 \in \c{A} \cap \c{B}$ be such that $\supp_A(v_0) \cap \supp_A(v_1) \neq \emptyset$. We may write:
        \begin{align}
            v_0 &= \sum_{i \in \supp_A(v_0)} \xi_i^0x_i = \sum_{j \in \supp_B(v_0)} \mu_j^0y_j, \label{eq:v_0} \\
            v_1 &= \sum_{i \in \supp_A(v_1)} \xi_i^1x_i = \sum_{j \in \supp_B(v_1)} \mu_j^1y_j. \label{eq:v_1}
        \end{align}
        where $\xi_i^0,\xi_i^1,\xi_j^0,\xi_j^1$ are all non-zero. We let:
        \begin{align*}
            w := \sum_{i \in \supp_A(v_0) \cap \supp_A(v_1)} \xi_i^0 x_i, \\
            w' := \sum_{j \in \supp_B(v_0) \cap \supp_B(v_1)} \mu_j^0 y_j.
        \end{align*}
        It's clear that $\supp_A(w) = \supp_A(v_0) \cap \supp_A(v_1)$. We shall show that $w = w'$. This implies that that $w \in \c{A} \cap \c{B}$, proving the claim.

        By the symmetry of the argument, it suffices to show that $\supp(w) \subseteq \supp(w')$, and $\lambda_k(w) = \lambda_k(w')$ for all $k \in \supp(w)$. Let $k \in \supp(w)$, so $k \in \supp(x_i)$ for some $i \in \supp_A(v_0) \cap \supp_A(v_1)$. By the second equalities in (\ref{eq:v_0}) and (\ref{eq:v_1}), we have that $k \in \supp(y_j)$ for some $j \in \supp_B(v_0)$, and $k \in \supp(y_{j'})$ for some $j' \in \supp_B(v_1)$. Since $B$ is a block sequence, $k \in \supp(y_j) \cap \supp(y_{j'})$ implies that $j = j'$. Therefore, $j \in \supp_B(v_0) \cap \supp_B(v_1)$, so $k \in \supp(w')$. By (\ref{eq:v_0}), we can conclude that:
        \begin{align*}
            \lambda_k(w) = \xi_i^0\lambda_k(x_i) = \lambda_k(v_0) = \mu_j^0\lambda_k(y_j) = \lambda_k(w'),
        \end{align*}
        which completes the proof. 
    \end{midproof}
    
    Let $K \subseteq \omega$ be the set of all $n$ such that $n \in \supp_A(v)$ for some $v \in \c{A} \cap \c{B}$. Since $\c{A} \cap \c{B}$ is infinite-dimensional, $K$ is infinite. Define a relation $\sim$ on $K$ such that for all $m,n \in K$:
    \begin{align*}
        m \sim n \iff \forall v \in \c{A} \cap \c{B}(m \in \supp_A(v) \leftrightarrow n \in \supp_A(v)).
    \end{align*}
    It's clear from the definition that $\sim$ is an equivalence relation, and that every equivalence class is finite. 

    \begin{claim}
        For each $n \in K$, there exists some $v \in \c{A} \cap \c{B}$ such that $\supp_A(v) = [n]_\sim$. 
    \end{claim}

    \begin{midproof}
        It follows from the definition of $\sim$ that:
        \begin{align*}
            [n]_\sim = \bigcap\{\supp_A(u) : u \in \c{A} \cap \c{B} \text{ and } n \in \supp_A(u)\}.
        \end{align*}
        However, since $\supp_A(u)$ is finite for each $u \in \c{A} \cap \c{B}$, we may find $u_0,\dots,u_{m-1} \in \c{A} \cap \c{B}$ such that $[n]_\sim = \bigcap_{i<m} \supp_A(u_i)$. Then the existence of such a $v$ follows from Claim \ref{claim:supp(v).intersection}.
    \end{midproof}

    Thus, for each $n$ we let $z_n \in \c{A} \cap \c{B}$ be a fixed vector such that $\supp_A(z_n) = [n]_\sim$. 

    \begin{claim}
    \label{claim:eq.classes.are.blocks}
        Let $m,n,l \in K$ be such that $m < n < l$, and suppose that $m \sim l$. Then $m \sim n \sim l$.
    \end{claim}

    \begin{midproof}
        Suppose otherwise, so there exist $m < n < l$ in $K$ such that $m \sim l$ but $n \notin [m]_\sim$. We write:
        \begin{align*}
            z_m &= \sum_{i \in \supp_A(z_m)} \xi_ix_i = \sum_{j \in \supp_B(z_m)} \mu_jy_j,
        \end{align*} 
        where $\xi_i,\mu_j$ are non-zero. Let $i_0 := \max\{i \in \supp_A(z_m) : i < n\}$ and $i_1 := \min\{i \in \supp_A(z_m) : i > n\}$, and let $k_0 := \max(\supp(x_{i_0}))$ and $k_1 := \min(\supp(x_{i_1}))$. We define:
        \begin{align*}
            w := \sum_{\substack{i \in \supp_A(z_m) \\ \max(\supp(x_i)) \leq k_0}} \xi_ix_i.
        \end{align*}
        We shall show that $w \in \c{A} \cap \c{B}$. This gives us the required contradiction, as $m \in \supp_A(w)$ but $l \not\in \supp_A(w)$. 

        Observe that for each $j \in \supp_B(z_m)$, either $k_0 < \min(\supp(y_j))$ or $\max(\supp(y_j)) < k_1$ --- otherwise, if $j \in \supp_B(z_m)$ is such that $\min(\supp(y_j)) \leq k_0 < k_1 \leq \max(\supp(y_j))$, then by the block sequence property of $B$, we have that for all $j' \neq j$, either $\max(\supp(y_{j'})) < k_0$ or $k_1 < \min(\supp(y_{j'}))$. Since $k_0 < \min(\supp(x_n)) \leq \max(\supp(x_n)) < k_1$ and $j \notin \supp_B(z_n)$, this implies that:
        \begin{align*}
            \supp(x_n) \cap \bigcup_{j' \in \supp_B(z_n)} \supp(y_{j'}) = \emptyset,
        \end{align*}
        which contradicts that $\supp(x_n) \subseteq \supp(z_n)$. 

        Note that for all $k \in \supp(z_m)$, either $k \leq k_0$ or $k_1 \leq k$. Thus, for all $j \in \supp_B(z_m)$, since $\supp(y_j) \subseteq \supp(z_m)$, if $\max(\supp(y_j)) < k_1$ then $\max(\supp(y_j)) \leq k_0$. Similarly, if $k_0 < \min(\supp(y_j))$ then $k_1 \leq \min(\supp(y_j))$. Therefore, we have that:
        \begin{align*}
            z_m = \sum_{\substack{j \in \supp_B(z_m) \\ \max(\supp(y_i)) \leq k_0}} \mu_j y_j + \sum_{\substack{j \in \supp_B(z_m) \\ k_1 \leq \min(\supp(y_i))}} \mu_j y_j.
        \end{align*}
        We may now express $\supp(z_m)$ in two different ways:
        \begin{align*}
            \bigcup_{i \in \supp_A(z_m)} \supp(x_i) &= \bigcup_{\substack{i \in \supp_A(z_m) \\ \max(\supp(x_i)) \leq k_0}} \supp(x_i) \cup \bigcup_{\substack{i \in \supp_A(z_m) \\ k_1 \leq \min(\supp(x_i))}} \supp(x_i), \\
            \bigcup_{j \in \supp_B(z_m)} \supp(y_j) &= \bigcup_{\substack{j \in \supp_B(z_m) \\ \max(\supp(y_j)) \leq k_0}} \supp(y_j) \cup \bigcup_{\substack{j \in \supp_B(z_m) \\ k_1 \leq \min(\supp(y_j))}} \supp(y_j).
        \end{align*}
        This implies that:
        \begin{align*}
            \bigcup_{\substack{i \in \supp_A(z_m) \\ \max(\supp(x_i)) \leq k_0}} \supp(x_i) &= \bigcup_{\substack{j \in \supp_B(z_m) \\ \max(\supp(y_j)) \leq k_0}} \supp(y_j), \\
            \bigcup_{\substack{i \in \supp_A(z_m) \\ k_1 \leq \min(\supp(x_i))}} \supp(x_i) &= \bigcup_{\substack{j \in \supp_B(z_m) \\ k_1 \leq \min(\supp(y_j))}} \supp(y_j).
        \end{align*}
        Therefore, if we let:
        \begin{align*}
            w' := \sum_{\substack{j \in \supp_B(z_m) \\ \max(\supp(y_j)) \leq k_0}} \mu_jy_j,
        \end{align*}
        then we have that $\supp(w) = \supp(w')$. Furthermore, for any $k \in \supp(w)$, $\lambda_k(w) = \lambda_k(z_m) = \lambda_k(w')$, so $w = w'$. It's clear from the definition that $w \in \c{A}$ and $w' \in \c{B}$, so we have that $w \in \c{A} \cap \c{B}$, as desired.
    \end{midproof}

    We now choose an increasing sequence $(n_i)_{i<\omega}$ of representatives of all equivalence classes in $\sim$ (i.e. $\bigcup_{i<\omega} [n_i]_\sim = K$). We let $C := (z_{n_i})_{i<\omega}$. By Claim \ref{claim:eq.classes.are.blocks}, $C$ is an infinite block sequence. We may finish the proof of this lemma by proving the following claim:

    \begin{claim}
        $\c{C} = \c{A} \cap \c{B}$.
    \end{claim}

    \begin{midproof}
        Since $z_{n_i} \in \c{A} \cap \c{B}$ for all $i$, we have that $\c{C} \subseteq \c{A} \cap \c{B}$. Conversely, let $v \in \c{A} \cap \c{B}$. Let $n_{i_0},\dots,n_{i_{t-1}}$ be an increasing enumeration of all $n_{i_j}$'s such that $n_{i_j} \in \supp_A(v)$. Fix any $k_j \in \supp(x_{n_{i_j}}) \subseteq \supp(z_{n_{i_j}})$, and let:
        \begin{align*}
            \xi_j := \frac{\lambda_{k_j}(v)}{\lambda_{k_j}(z_{n_{i_j}})}.
        \end{align*}
        Note that $\xi_j \neq 0$ for all $j$. Now let:
        \begin{align*}
            v' := \sum_{j<t} \xi_jz_{n_{i_j}}.
        \end{align*}
        Clearly $v' \in \c{C}$. We shall show that $v = v'$, which proves the claim.

        We first note that $\supp_A(v') = \supp_A(v)$ (which implies that $\supp(v') = \supp(v)$) --- observe that $\supp_A(v') = \bigcup_{j<t} \supp_A(z_{n_{i_j}})$. If $m \in \supp_A(v)$, then $[m]_\sim \subseteq \supp_A(v)$ and $[m]_\sim = [n_{i_j}]_\sim = \supp_A(z_{n_{i_j}})$ for some $j$, so $m \in \supp_A(v')$. Conversely, if $m \in \supp_A(z_{n_{i_j}})$, then since $n_{i_j} \in \supp_A(v)$, we have that $[m]_\sim = [n_{i_j}]_\sim \subseteq \supp_A(v)$, so $m \in \supp_A(v)$. 

        It remains to show that $\lambda_k(v') = \lambda_k(v)$ for all $k \in \supp(v)$. We fix some $j < t$, and observe that:
        \begin{align*}
            \lambda_{k_j}(v') = \frac{\lambda_{k_j}(v)}{\lambda_{k_j}(z_{n_{i_j}})} \cdot \lambda_{k_j}(z_{n_{i_j}}) = \lambda_{k_j}(v).
        \end{align*}
        Since $v \in \c{A}$ and $k_j \in \supp(x_{n_{i_j}}) \subseteq \supp(v)$, there exist some vectors $u,w \in \c{A}$ (or $u = 0$ or $w = 0$) such that $u < x_{n_{i_j}} < w$, and:
        \begin{align*}
            v = u + \frac{\lambda_{k_j}(v)}{\lambda_{k_j}(x_{n_{i_j}})}x_{n_{i_j}} + w.
        \end{align*}
        Similarly, there exist some vectors $u',w' \in \c{A}$ such that $u' < x_{n_{i_j}} < w'$ (or $u' = 0$ or $w'      = 0$), and:
        \begin{align*}
            v' = u' + \frac{\lambda_{k_j}(v)}{\lambda_{k_j}(x_{n_{i_j}})}x_{n_{i_j}} + w'.
        \end{align*}
        Therefore, $v - v' \in \c{A} \cap \c{B}$, and $n_{i_j} \notin \supp(v - v')$. But this implies that $[n_{i_j}]_\sim \cap \supp(v - v') = \emptyset$, which is only possible if for all $k \in \supp(z_{n_{i_j}})$, $\lambda_k(v') = \lambda_k(v)$. Therefore, for all $j < t$ and $k \in \supp(z_{n_{i_j}})$, $\lambda_k(v') = \lambda_k(v)$, so $v' = v$, as desired.
    \end{midproof}
\end{proof}

\end{document}
